\zhangtitle{第1章 空间向量与立体几何}
\statsline{全件140题：简单21｜中档104｜难15}
\vspace{4pt}
\noindent
\begin{tabular}{|>{\raggedright\arraybackslash}p{78mm}|>{\centering\arraybackslash}p{14mm}|>{\raggedright\arraybackslash}p{44mm}|>{\centering\arraybackslash}p{30mm}|}
\hline\multicolumn{1}{|>{\centering\arraybackslash}p{78mm}|}{\textbf{节名}} & \multicolumn{1}{|>{\centering\arraybackslash}p{14mm}|}{\textbf{题量}} & \multicolumn{1}{|>{\centering\arraybackslash}p{44mm}|}{\textbf{简单/中档/难}} & \multicolumn{1}{|>{\centering\arraybackslash}p{30mm}|}{\textbf{题型组数}} \\ \hline
1.1.1 空间向量及其运算 & 10 & 简单6/中档4/难0 & 9 \\ \hline
1.1.2 空间向量基本定理 & 4 & 简单2/中档2/难0 & 3 \\ \hline
1.1.3 空间向量的坐标与空间直角坐标系 & 10 & 简单3/中档7/难0 & 7 \\ \hline
1.2.1 空间中的点、直线与空间向量 & 9 & 简单2/中档6/难1 & 6 \\ \hline
1.2.2 空间中的平面与空间向量 & 7 & 简单2/中档5/难0 & 5 \\ \hline
1.2.3 直线与平面的夹角 & 8 & 简单1/中档6/难1 & 7 \\ \hline
1.2.4 二面角 & 13 & 简单0/中档13/难0 & 12 \\ \hline
1.2.5 空间中的距离 & 79 & 简单5/中档61/难13 & 69 \\ \hline
合计 & 140 & 简单21/中档104/难15 & 118 \\ \hline
\end{tabular}
\vspace{14pt}
\jietitle{1.1 空间向量及其运算}
\xiaojietitle{1.1.1 空间向量及其运算}
\statsline{本节10题：简单6｜中档4｜难0　提分线：简单→保60\%｜中档→保80\%}
